\section{整体灵敏度与鲁棒性分析}
\label{sec:overall-sensitivity}

\subsection{核数变化下的速度比}
主矩阵对100张图、1--5核和A、B、C三场景各自独立生成候选并调用官方评价器。A场景平均速度比由1核的1.0220升至五核3.4383，B由1.0564升至3.4738，C由1.0739升至3.5533。核数增加带来的收益来自可并行Task增多，同时受到依赖最长路、跨核释放、片上容量和共享DDR服务的限制。

\begin{table}[H]\centering\small
\caption{三场景逐核平均速度比（v7主矩阵）}\label{tab:overall-speedup-v7}
\begin{tabular}{c rrr}\toprule
核数 & A场景 & B场景 & C场景\\\midrule
1 & 1.0220 & 1.0564 & 1.0739\\
2 & 1.7432 & 1.7610 & 1.7998\\
3 & 2.3742 & 2.3993 & 2.4462\\
4 & 2.9520 & 2.9852 & 3.0448\\
5 & 3.4383 & 3.4738 & 3.5533\\\bottomrule
\end{tabular}
\end{table}

对每个组合，最终工期、计划哈希和候选来源均由该组合的\texttt{final\_selection}确定；正文不再将低核计划事后复制为高核结果。由有放回重采样得到的均值区间随图集合而变，图\ref{fig:speedup-a}--\ref{fig:speedup-c}的浅色带用于表示该实例组成敏感性。

\subsection{候选选择消融}
初始方案与最终方案使用同一候选台账和同一官方评价器，差异只来自有限候选的排序、评价和选优。五核平均改善分别为A的58.6697\%（89/100）、B的60.3622\%（98/100）和C的1.1744\%（38/100）。A、B的改善较大，原因是分组切分、组移动和边界变体覆盖了更多有效的并行及容量折中；C场景从B计划起步，Cache候选的可行改动空间较窄。

\begin{table}[H]\centering\small
\caption{五核初始方案到最终方案的候选选择收益}\label{tab:overall-ablation-v7}
\begin{tabular}{lrrr}\toprule
场景 & 平均工期改善/\% & 改善图数 & 官方成功组合\\\midrule
A & 58.6697 & 89/100 & 100\\
B & 60.3622 & 98/100 & 100\\
C & 1.1744 & 38/100 & 100\\\bottomrule
\end{tabular}
\end{table}

\subsection{Cache三段效应}
500组B/C同计划配对逐组核对计划身份，并计算硬件效应、C内调度适配和综合效应。五核均值分别为1.0146、1.0139和1.0289；三者逐组满足乘积恒等式，列均值只作为等权汇总。完整结果与公式见第\ref{sec:question-three}节。

\subsection{Cache容量与带宽敏感性}
六张代表图分别进行容量减半、容量加倍、带宽减半和带宽加倍。24组固定计划全部通过官方评价，37条重优化候选全部成功，每个变体均有3/6组改变计划。重优化相对默认工期比均值依次为0.9900、0.9869、0.9892和0.9885；容量加倍的下降幅度最大。固定计划下的尾部工期在这24组中保持默认值，命中率变化主要在重优化后通过不同候选的事件时序体现。

\begin{table}[H]\centering\small
\caption{Cache硬件扰动的代表图实验}\label{tab:overall-cache-sensitivity}
\begin{tabular}{lrrrr}\toprule
扰动 & 固定命中率 & 重优化命中率 & 重优化/默认工期比 & 改变计划\\\midrule
容量减半 & 2.0977\% & 3.4413\% & 0.9900 & 3/6\\
容量加倍 & 8.1117\% & 13.0176\% & 0.9869 & 3/6\\
带宽减半 & 4.7811\% & 7.2801\% & 0.9892 & 3/6\\
带宽加倍 & 4.7811\% & 7.3484\% & 0.9885 & 3/6\\\bottomrule
\end{tabular}
\end{table}

命中率和工期不构成一一对应关系：只有当读请求位于决定尾部完成时刻的路径上，Cache服务来源变化才会改变最终Makespan。因而参数分析同时报告命中率、候选计划是否改变和重优化后的官方工期。

\subsection{结果边界}
v7主矩阵是有限候选启发式在题给CPU Python官方事件模拟器上的结果，不能表述为全局最优证明，也不能表述为NPU实机测量。主矩阵中228个组合由fast-profile尾部补齐，仍通过官方评价并保留来源；支撑记录另行登记solver文件哈希漂移，以及case\_054/5核、case\_062/4核和case\_085/5核三个B/C初始计划不一致项。上述边界只影响证据解释，不改变论文使用的最终选择文件和汇总口径。
